\chapter{Beyond Almgren--Chriss}\label{ch:mx:beyond-almgren-chriss}

The Almgren--Chriss schedule does not look at the market while it trades. A trader who expects the price to fall in the next half-hour should sell
faster; one whose impact decays should pause after a large trade; one whose market dries up at lunch should trade before it does. This chapter adds
each of these to chapter~14's block: a forecast, the price's own moves, transient impact and random liquidity, and measures what each is worth.

\section{Execution with a signal}

\begin{definition}[Signal-adaptive schedule]\label{def:mx:beyond-almgren-chriss:signal}
A \emph{signal-adaptive schedule}\index{signal-adaptive schedule} sets each interval's trade from the holdings left and a forecast of the price's drift,
trading faster when the forecast is against the order (a seller expecting a fall) and slower when it is for it.
\end{definition}

Let the price drift at a rate $\alpha_t$ that follows an Ornstein--Uhlenbeck process (One Quant Book 4, chapter~4) with a half-life and a stationary
standard deviation, observed by the trader. On $n$ steps of length $\tau$, selling $u_k$ from holdings $x_k$, the expected cost is quadratic in the state
$(x_k,\alpha_k)$ and the control:
\[
\min\ \E\sum_k\Bigl(\eta\frac{u_k^2}{\tau}+\lambda\sigma^2\tau x_k^2-\tau\,\alpha_kx_k\Bigr),\qquad x_{k+1}=x_k-u_k,\quad \alpha_{k+1}=\phi\,\alpha_k+\text{noise},
\]
with $x_n=0$ enforced by a large terminal penalty. A linear-quadratic problem has a linear optimal policy, $u_k=K_{k,1}x_k+K_{k,2}\alpha_k$, with gains from a
backward Riccati recursion; without a signal it reproduces the Almgren--Chriss trajectory exactly. Cartea and Jaimungal (2016) built order flow into the
same framework.

\omcode{code/firm/execcontrol/firm_execcontrol.py}{35}{47}{The Riccati recursion of execution with a signal: gains on the holdings and on the forecast drift.}\label{lst:mx:beyond-almgren-chriss:lqr}

On the block of chapter~14 (1 million shares of a USD~50 stock over a day, in 78 five-minute steps, $\lambda=10^{-6}$), 1\,500 simulated days with the same
random numbers for both schedules give the saving in \cref{fig:mx:beyond-almgren-chriss:signal}. A signal with a half-life of a twentieth of a day (about
twenty minutes) and a standard deviation of USD~0.5 a day saves 0.19 basis points against Almgren--Chriss's 61; one with a half-life of a day and
USD~2 a day saves 41 of 66. The saving grows with the signal's strength and more than linearly with its persistence: a forecast that lasts only a few
steps cannot be acted on before it is gone. It also raises the cost's variance (a standard deviation of 212 basis points against 173 for the strongest
signal): the schedule takes risk to use the forecast.

\begin{omfigure}
\begin{tikzpicture}
\begin{axis}[width=9cm,height=5.4cm,xmode=log,xlabel={\small signal half-life (days)},ylabel={\small saving (basis points)},
  tick label style={font=\footnotesize},xtick={0.05,0.2,1},xticklabels={0.05,0.2,1},grid=major,grid style={gray!25},ymin=0,ymax=45,
  legend columns=3,legend style={font=\footnotesize,at={(0.5,-0.3)},anchor=north,draw=none,fill=none,
  /tikz/every even column/.append style={column sep=8pt}},table/col sep=comma]
\addplot[omDef,very thick,mark=*,error bars/.cd,y dir=both,y explicit] table[x=half,y=s0,y error=e0]{figdata/microstructure/15-beyond-almgren-chriss/signal.csv};
\addplot[omThm,very thick,mark=square*,error bars/.cd,y dir=both,y explicit] table[x=half,y=s1,y error=e1]{figdata/microstructure/15-beyond-almgren-chriss/signal.csv};
\addplot[omProp,very thick,mark=triangle*,error bars/.cd,y dir=both,y explicit] table[x=half,y=s2,y error=e2]{figdata/microstructure/15-beyond-almgren-chriss/signal.csv};
\legend{USD 0.5 a day,USD 1 a day,USD 2 a day}
\end{axis}
\end{tikzpicture}

\omcaption{What a drift forecast saves the block's execution, against its half-life, for three strengths (standard deviations of the drift): mean saving
over 1\,500 simulated days with common random numbers, and its standard error. Data: \texttt{mx\_beyond.signal\_study}.}\label{fig:mx:beyond-almgren-chriss:signal}
\end{omfigure}

\section{Adaptive schedules}

\begin{definition}[Adaptive execution, aggressive-in-the-money strategy]\label{def:mx:beyond-almgren-chriss:adaptive}
An \emph{adaptive execution}\index{adaptive execution} changes its schedule with the price's path during the order. An \emph{aggressive-in-the-money
strategy}\index{aggressive-in-the-money strategy} trades faster when the price has moved in the order's favour (up for a seller) and slower when it has
moved against it.
\end{definition}

Almgren--Chriss is static: its mean--variance optimum is computed once. Lorenz and Almgren (2011) showed that re-optimising as the price moves does better
in the same mean--variance sense. The intuition: after a favourable move the order has gained money it can spend on impact to cut the remaining risk.
The chapter's rule recomputes the Almgren--Chriss trajectory from the current holdings at each step with \omterm{def:mx:the-almgren-chriss-framework:urgency}{urgency} $\kappa_0e^{a\,\Delta p/(\sigma\sqrt T)}$.

\omcode{code/firm/execcontrol/firm_execcontrol.py}{78}{91}{An aggressive-in-the-money policy: the static trajectory re-solved from the current holdings with an urgency that grows with the gain so far.}\label{lst:mx:beyond-almgren-chriss:aim}

With $\kappa_0T=2$ and 3\,000 simulated days, the static schedule costs 60.9 basis points with a standard deviation of 90.5, a certainty equivalent of 101.9. The
adaptive rule with $a=0.25$, 0.5, 1 and 2 costs 61.2, 62.2, 66.2 and 81.7 on average with standard deviations of 88.7, 86.4, 80.4 and 66.2, and certainty
equivalents of 100.5, 99.4, 98.6 and 103.6: moderate adaptation improves the trade-off by up to 3.2\%, strong adaptation overshoots.

\begin{omfigure}
\begin{tikzpicture}
\begin{axis}[width=8.5cm,height=5.2cm,xlabel={\small standard deviation (basis points)},ylabel={\small mean cost (basis points)},
  tick label style={font=\footnotesize},xmin=60,xmax=95,ymin=55,ymax=88,grid=major,grid style={gray!25},table/col sep=comma]
\addplot[omDef,very thick,mark=*] table[x=sd,y=mean]{figdata/microstructure/15-beyond-almgren-chriss/aim.csv};
\node[font=\scriptsize,anchor=south west] at (axis cs:90.5,60.9) {static};
\node[font=\scriptsize,anchor=west] at (axis cs:80.4,66.2) {$a=1$};
\node[font=\scriptsize,anchor=west] at (axis cs:66.2,81.7) {$a=2$};
\end{axis}
\end{tikzpicture}

\omcaption{The aggressive-in-the-money rule for increasing adaptation $a=0$ (static), 0.25, 0.5, 1 and 2: it trades mean cost for lower risk, with the best
certainty equivalent near $a=1$ at this risk aversion. 3\,000 simulated days. Data: \texttt{mx\_beyond.aim\_grid}.}\label{fig:mx:beyond-almgren-chriss:aim}
\end{omfigure}

\section{Transient impact and the optimal block}

\begin{definition}[Bucket-shaped trajectory]\label{def:mx:beyond-almgren-chriss:bucket}
A \emph{bucket-shaped trajectory}\index{bucket-shaped trajectory} trades a block at the start and another at the end with a slower constant rate between:
the optimum when impact is transient.
\end{definition}

With the exponential kernel of Obizhaeva and Wang (2013, chapter~12), resilience $\rho=10$ per day and 101 trades over the day, the time-weighted schedule
costs 0.0893 (in units of the instantaneous impact of the whole order) and the optimum 0.0834, 6.6\% less: it trades 0.0876 of the order in each end block,
close to the formula's $1/(\rho T+2)=0.0833$, and 0.0083 in each trade between. Selling everything at once would cost 0.5. The saving is modest because the
kernel decays fast relative to the day; with slower resilience the blocks and the saving grow.

\section{Stochastic liquidity}

\begin{definition}[Stochastic liquidity]\label{def:mx:beyond-almgren-chriss:stochliq}
\emph{Stochastic liquidity}\index{stochastic liquidity} lets the impact parameters themselves change at random during the execution, as volume and
depth do, so that the cost of trading now depends on the state of the market now.
\end{definition}

Almgren (2012) solved mean--variance execution when liquidity and volatility vary randomly, with strategies that adapt to the market's quality. The
chapter's version has two regimes, normal and dry, the dry one with four times the temporary impact; the market stays normal with probability 0.9 per
step and dry with 0.7; the policy on (holdings, regime) comes from Book~4's \texttt{firm.dpsolve}. Halfway through the day with half the order left, it
sells 60\,000 shares in the next step when the market is normal and 20\,000 when it is dry. Over 3\,000 simulated days it costs 73.2 basis points against
91.5 for Almgren--Chriss, which ignores the regimes: 18.3 basis points saved (standard error 0.5), with a slightly lower standard deviation (89.6 against 94.6).

\section{A control-theory template}

The four problems share a template: a state (holdings, and whatever the market reveals), a control (the trade), a cost (impact now, risk and forecast
drift later), and a solution method matched to the structure: a closed form for Almgren--Chriss, a Riccati recursion for linear-quadratic signals, a
matrix inverse for transient impact, dynamic programming for regimes. Gatheral and Schied (2011) solved the geometric Brownian case in closed form under
another risk criterion. Every solution returns a schedule in \texttt{firm.acexec}'s interface, so the child-order layer does not care which one produced it.

\section{Tutorial: trading on a forecast}\label{tut:mx:beyond-almgren-chriss:tut}

\begin{tutorial}
\textbf{Goal.} Add a signal, adaptation, transient impact and random liquidity to the block's execution, and price each.
\textbf{End state:} \cref{fig:mx:beyond-almgren-chriss:signal,fig:mx:beyond-almgren-chriss:aim} and the numbers of sections 2 to 4.
\begin{tutsteps}
\item \textbf{Signal.} \texttt{firm\_execcontrol.lqr\_signal(n, tau, eta, sigma, lam, phi)} and \texttt{simulate(policy, \dots)}; \texttt{mx\_beyond.signal\_study()}.
\item \textbf{Adaptation.} \texttt{aim(k0, a, sigma, T, n)}; \texttt{aim\_grid()}.
\item \textbf{Transient.} \texttt{transient\_study(rho)} with \texttt{firm\_propagator.optimal\_liquidation} and \texttt{OWScheduler}.
\item \textbf{Liquidity.} \texttt{liquidity\_dp(X, n, tau, eta, sigma, lam, stay)} and \texttt{liquidity\_study()}; draw with \texttt{fig\_beyond.py}.
\end{tutsteps}
\textbf{What to change next.} Feed each schedule to chapter~14's \texttt{ScheduleAgent} and run it in the simulated market on many sessions; estimate the
signal's drift with a Kalman filter instead of observing it.
\end{tutorial}

\section{Build: execution control}\label{bld:mx:beyond-almgren-chriss:execcontrol}

\begin{build}
\bfield{Purpose} The adaptive schedules behind the algorithms of chapters 16 and 28 and Book~12's learned execution policies, as benchmarks with known
optima.
\bfield{Interface} \texttt{lqr\_signal(n, tau, eta, sigma, lam, phi)}, \texttt{simulate(policy, X, n, tau, eta, sigma, phi, s\_alpha, seed, paths)},
\texttt{OWScheduler(X, rho, T)}, \texttt{aim(k0, a, sigma, T, n)}, \texttt{liquidity\_dp(X, n, tau, eta, sigma, lam, stay)}.
\bfield{Rules} Sales are positive; the last step sells what is left; costs against the arrival price per share; common random numbers across policies;
regime probabilities in tenths (the nodes of \texttt{firm.dpsolve}).
\bfield{Acceptance tests} \texttt{code/firm/execcontrol/tests/}: the Riccati policy without a signal equals the discrete Almgren--Chriss trajectory; a
signal lowers the mean cost; Obizhaeva--Wang's blocks; the adaptive policy's direction and its static limit; more selling when liquid.
\bfield{Stretch} Kalman-filtered signals; stochastic volatility; transient impact with a signal.
\end{build}

\begin{omsources}
\item P.~Lorenz and R.~Almgren, ``Mean--variance optimal adaptive execution'', \emph{Applied Mathematical Finance} 18(5), 2011.
\item J.~Gatheral and A.~Schied, ``Optimal trade execution under geometric Brownian motion in the Almgren and Chriss framework'', \emph{International
Journal of Theoretical and Applied Finance} 14(3), 2011.
\item R.~Almgren, ``Optimal trading with stochastic liquidity and volatility'', \emph{SIAM Journal on Financial Mathematics} 3(1), 2012.
\item A.~A.~Obizhaeva and J.~Wang, ``Optimal trading strategy and supply/demand dynamics'', \emph{Journal of Financial Markets} 16(1), 2013.
\item A.~Cartea and S.~Jaimungal, ``Incorporating order-flow into optimal execution'', \emph{Mathematics and Financial Economics} 10(3), 2016.
\end{omsources}

\section{Exercises}

\begin{exercise}[$\star$]\label{exo:mx:beyond-almgren-chriss:1}
A signal has a half-life of 0.2 days. What is its persistence $\phi$ over a five-minute step of a 78-step day?
\end{exercise}

\begin{exercise}[$\star$]\label{exo:mx:beyond-almgren-chriss:2}
Why is the gain on the forecast drift negative in the Riccati policy for a seller (the policy sells less when the forecast drift is positive)?
\end{exercise}

\begin{exercise}[$\star$]\label{exo:mx:beyond-almgren-chriss:3}
With resilience $\rho=10$ per day and a one-day window, what share of the order does the Obizhaeva--Wang optimum trade in each end block?
\end{exercise}

\begin{exercise}[$\star\star$]\label{exo:mx:beyond-almgren-chriss:4}
Why does a signal with a short half-life save so little, even when it is strong?
\end{exercise}

\begin{exercise}[$\star\star$]\label{exo:mx:beyond-almgren-chriss:5}
Why does adaptation with $a=2$ do worse than the static schedule in certainty equivalent?
\end{exercise}

\begin{exercise}[$\star\star$]\label{exo:mx:beyond-almgren-chriss:6}
What does the stochastic-liquidity policy do differently from Almgren--Chriss, and why does it save 18 basis points?
\end{exercise}

\begin{exercise}[$\star\star\star$]\label{exo:mx:beyond-almgren-chriss:7}
\emph{Coding.} Run the signal study with the signal observed with noise (the policy sees $\alpha+\epsilon$ with $\epsilon$ of the same standard deviation as
$\alpha$). How much of the saving survives, and why?
\end{exercise}

\begin{exercise}[$\star\star\star$]\label{exo:mx:beyond-almgren-chriss:8}
\emph{Find the flaw.} ``Our execution with the alpha signal saved 40 basis points in the backtest, so the signal is worth 40 basis points on every order.''
\end{exercise}

\section{Problem: Trading on a Forecast}

\begin{problem}[{Weekend problem --- trading on a forecast}]\label{pb:mx:beyond-almgren-chriss:1}
A research team gives the execution desk a drift forecast and asks what it is worth; the desk also wonders about adapting to prices and liquidity.

\medskip\noindent\textbf{Part I --- The signal.}
\begin{enumerate}
\item Write the linear-quadratic problem and its state.
\item Why is the optimal policy linear, and how are its gains computed?
\item Check that without a signal it is Almgren--Chriss.
\item Give the saving for the three strengths at the three half-lives.
\end{enumerate}

\medskip\noindent\textbf{Part II --- Adaptation.}
\begin{enumerate}[resume]
\item Define an \omterm{def:mx:beyond-almgren-chriss:adaptive}{aggressive-in-the-money strategy}.
\item Give the static schedule's cost, risk and certainty equivalent.
\item Give the adaptive rule's for $a=0.25$, 0.5, 1 and 2.
\item Why does moderate adaptation help?
\end{enumerate}

\medskip\noindent\textbf{Part III --- Transient impact and liquidity.}
\begin{enumerate}[resume]
\item Give the costs of the time-weighted and optimal schedules under the transient kernel and the size of the blocks.
\item Why is the saving small here?
\item Describe the two liquidity regimes and the policy's trades in each.
\item Give the dynamic programme's cost and saving.
\end{enumerate}

\medskip\noindent\textbf{Part IV --- The verdict.}
\begin{enumerate}[resume]
\item Rank the four extensions by what they saved here.
\item Which of them need data the desk may not have?
\item What does the signal do to the cost's risk?
\item How would you test the signal's value on real orders?
\item State the \emph{named result}: the basis points an alpha signal saves an execution schedule, against the signal's half-life and strength.
\item What minimum half-life makes a signal worth wiring into execution, on these numbers?
\item Which extension would you build first, and why?
\item In one sentence: what does looking at the market buy an execution?
\end{enumerate}
\end{problem}

\section{Interview questions}

\begin{interviewq}[$\star$ \iqroles{trader}]\label{iq:mx:beyond-almgren-chriss:1}
You are selling and your alpha model says the stock will fall over the next hour. How does your schedule change?
\end{interviewq}

\begin{interviewq}[$\star\star$ \iqroles{researcher}]\label{iq:mx:beyond-almgren-chriss:2}
Set up execution with an Ornstein--Uhlenbeck signal as a linear-quadratic control problem.
\end{interviewq}

\begin{interviewq}[$\star\star$ \iqroles{researcher}]\label{iq:mx:beyond-almgren-chriss:3}
Why is the optimal execution under transient impact bucket-shaped?
\end{interviewq}

\begin{interviewq}[$\star\star$ \iqroles{trader}]\label{iq:mx:beyond-almgren-chriss:4}
Should an execution speed up or slow down after the price has moved in its favour?
\end{interviewq}

\begin{interviewq}[$\star\star$ \iqroles{developer}]\label{iq:mx:beyond-almgren-chriss:5}
How would you compare two execution policies fairly in simulation?
\end{interviewq}

\begin{interviewq}[$\star\star\star$ \iqroles{researcher}]\label{iq:mx:beyond-almgren-chriss:6}
How would you estimate the value of a short-horizon signal for execution from historical orders?
\end{interviewq}
